\documentclass[conference]{IEEEtran}
\IEEEoverridecommandlockouts

\usepackage{cite}
\usepackage{amsmath,amssymb,amsfonts}
\usepackage{graphicx}
\usepackage{textcomp}
\usepackage{xcolor}

\def\BibTeX{{\rm B\kern-.05em{\sc i\kern-.025em b}\kern-.08em
    T\kern-.1667em\lower.7ex\hbox{E}\kern-.125emX}}

\begin{document}

\title{Literature Review: An Artificial Neural Network-based Quantum Finance System}

\author{\IEEEauthorblockN{1\textsuperscript{st} Xiong Haoyu}
\IEEEauthorblockA{2330025090}
\and
\IEEEauthorblockN{2\textsuperscript{nd} Alexey Krivorotov}
\IEEEauthorblockA{2330026001}
\and
\IEEEauthorblockN{3\textsuperscript{rd} Chen Yiduo}
\IEEEauthorblockA{2330026025}
\and
\IEEEauthorblockN{4\textsuperscript{nd} Iaroslav Obukhov}
\IEEEauthorblockA{2330026065}
}

\maketitle

\section{Literature Review}

\subsection{Conceptual Overview of ANN}
We have reviewed a series of works and studies which talks about the use of Artificial Neural Networks, converting financial historical data and relevant features into forecasting. And the input is usually setting as a time-orded record like returns, prices and index values. The output differs by task: Wang et al. forecast next-day closing prices for SSE, TWSE, KOSPI, and Nikkei 225 indices \cite{wang2016elman}; Shen et al. classify index observations into long, short, and flat trading signals \cite{shen2018gru}; Sun et al. predict whether a firm's next half-year financial condition will be healthy or distressed \cite{sun2011sffs}. Therefore, the technology is better understood as neural financial sequence modeling than as simple price prediction.

These papers have covered some different sampling frequencies and markets that affects what the models can claim. Singh and Srivastava use daily Google OHLC data plus 36 technical variables before $(2D)^2$PCA and DNN training \cite{singh2017stock}; Chong et al. use ten lagged 5-minute returns for 38 KOSPI stocks and compare raw, PCA, RBM, and autoencoder representations \cite{chong2017deep}; Bao et al. test six major stock-index markets with wavelet denoising, stacked autoencoders, and LSTM \cite{bao2017deep}. Trading-oriented studies move further toward decisions: Fischer and Krauss predict next-day S\&P 500 cross-sectional winners from 1992 to 2015 \cite{fischer2018deep}, Gunduz et al. and Lachiheb and Gouider study intraday equity direction \cite{gunduz2017intraday,lachiheb2018hierarchical}, Nakano et al. build Bitcoin trading positions from 15-minute data \cite{nakano2018bitcoin}, and Evans et al. combine ANN with genetic algorithms for intraday foreign exchange \cite{evans2013utilizing}. The common thread is nonlinear pattern learning, but the supported decisions range from forecasting to trading execution and distress warning.

\subsection{Methodological Mechanism}
The first step is representation. Singh and Srivastava turn daily Google prices and technical indicators into chart-like matrices, then apply $(2D)^2$PCA before the DNN; their best setting, a 20-day window and $10 \times 10$ reduced dimension, gives a 0.682 hit rate and 1.365 total return \cite{singh2017stock}. Chong et al. show why this step matters: DNNs outperform AR(10) clearly in training but only slightly in testing, while applying DNN to AR residuals reduces error more convincingly than the reverse experiment \cite{chong2017deep}. Bao et al. address market noise more directly with a WT-SAE-LSTM pipeline, using wavelet transform for denoising, stacked autoencoders for feature extraction, and LSTM for final price prediction \cite{bao2017deep}.

The second step is matching architecture to data structure. Recurrent models are used when temporal order is central. Wang et al.'s ST-ERNN forecasts next-day index closes from OHLC inputs and reports lower MAE, RMSE, and MAPE than BPNN, STNN, and ERNN baselines \cite{wang2016elman}. Fischer and Krauss use 240 standardized daily returns, approximately one trading year, as the LSTM input for S\&P 500 stocks \cite{fischer2018deep}. Shen et al. adopt a similar 240-day sequence and report that GRU-SVM reaches 51.7\% S\&P 500 accuracy, only slightly above GRU at 51.4\%, but clearly above DNN and SVM baselines \cite{shen2018gru}. These studies support recurrent backbones for market sequences, while also showing that classification improvements can be numerically modest.

Other papers emphasize structure across features or assets rather than only time. Gunduz et al. arrange hourly BIST features by hierarchical clustering of feature correlations; their correlation-ordered CNN reaches a mean macro F-measure of 0.563, above randomly ordered CNN features and logistic regression \cite{gunduz2017intraday}. Lachiheb and Gouider's hierarchical DNN uses both a target stock's past returns and other TUNINDEX stocks' returns, performing best for 60\% of the ten tested stocks under MSE and reaching up/down accuracy as high as 73\% \cite{lachiheb2018hierarchical}. Evaluation then becomes a financial question, not just a prediction question. Singh and Srivastava report that DNN is better on hit rate and return correlation, while RBFNN is better on total return and RMSE \cite{singh2017stock}; Fischer and Krauss therefore test a long-short portfolio and report 0.46\% average daily return before transaction costs and 0.26\% after costs \cite{fischer2018deep}; Nakano et al. similarly test Bitcoin strategies under 0, 2.5, 5, and 10 bps bid--ask spreads \cite{nakano2018bitcoin}. These disagreements make clear that a model should be judged as a decision rule, not only as a statistical predictor.

\subsection{Recent Research Advances}
The later studies move away from simple ANN forecasting and toward hybrid systems. Sun et al. combine feature selection, PCA-based financial situation evaluation, and GA-optimized BPNN for dynamic distress prediction; the important point is the rolling design that uses period $t-0.5$ indicators to predict period $t$ and then checks the result against ex-post PCA evaluation \cite{sun2011sffs}. Wang et al. similarly add a stochastic time effective function to recurrent learning and note that large market fluctuations increase relative errors, making short-horizon prediction more reliable than longer-horizon prediction \cite{wang2016elman}. These papers already treat financial prediction as time-varying rather than static fitting.

Contemporary deep financial prediction also relies heavily on preprocessing and representation learning. Chong et al. are cautious about DNN superiority because the test-set advantage over AR(10) is small, although the residual experiment suggests that DNNs capture information missed by a linear autoregressive model \cite{chong2017deep}. Singh and Srivastava show that window length and reduced dimension materially affect results, with oversized windows and high-dimensional settings weakening performance \cite{singh2017stock}. Bao et al. combine WT, SAE, and LSTM and evaluate both accuracy and profitability across six stock-index markets \cite{bao2017deep}. The direction of the literature is therefore not ``deep learning replaces feature engineering,'' but hybrid modeling that joins denoising, representation learning, and sequence models.

The most decision-oriented papers add a stricter warning. Fischer and Krauss give strong large-scale evidence for LSTM trading, but also show that after 2010 excess returns fluctuate around zero after transaction costs \cite{fischer2018deep}. Shen et al. improve GRU classification by replacing the softmax output with an SVM-style layer, yet the S\&P 500 accuracy remains only slightly above 50\% \cite{shen2018gru}. Nakano et al. find that Bitcoin ANN strategies can beat buy-and-hold under some cost assumptions, but fail when the spread is widened to 10 bps \cite{nakano2018bitcoin}. Evans et al. are also relevant because they connect ANN outputs to GA-based intraday FX trading design rather than stopping at prediction \cite{evans2013utilizing}. The resulting gap is practical: robustness across regimes, transaction costs, interpretability, portfolio-level rules, and whether the model improves financial outcomes beyond a narrow dataset.

The newest direction in the added literature is quantum computing for finance, which broadens the review from prediction to computation and optimization. Orus et al. explicitly frame quantum annealing as a tool to ``optimize portfolios, find arbitrage opportunities, and perform credit scoring,'' and also discuss QML for classification, regression, PCA, and neural-network training \cite{orus2019quantum}. Their review further links quantum amplitude estimation to a ``quantum speed-up for Monte Carlo sampling,'' which is relevant to derivative pricing, VaR, and CVaR \cite{orus2019quantum}. Herman et al. later organize the field more broadly around stochastic modeling, optimization, and machine learning, while stressing that end-to-end quantum advantage still faces data-loading, readout, and NISQ hardware limits \cite{herman2023quantum}. This does not replace the ANN literature above, but it gives a stronger basis for treating quantum finance as a computational layer for financial engineering problems.

\subsection{Integration with Quantum Finance and Applications in Financial Engineering}
The earlier ANN papers do not directly test quantum finance, but the added quantum-finance papers make the integration more concrete. Orus et al. argue that many financial problems have a natural optimization or Monte Carlo structure: portfolio optimization can be mapped to quantum annealing, optimal arbitrage can be written as a QUBO problem, credit-scoring feature selection can also be optimized by annealing, and quantum amplitude estimation can accelerate Monte Carlo estimation for derivative pricing and risk measures \cite{orus2019quantum}. They also discuss QML as a possible way to improve classification, regression, PCA, and neural-network training, although they caution that many QML algorithms still require more advanced universal quantum computers \cite{orus2019quantum}. Herman et al. strengthen the same point from a more recent perspective by separating stochastic modeling, optimization, and machine learning, and by noting that data loading, model readout, and error correction remain major barriers to practical quantum advantage \cite{herman2023quantum}.

A defensible integration is therefore to use neural networks to estimate predictive signals and then pass these signals into a quantum-inspired optimization or risk layer. The input design can follow the reviewed evidence: OHLC and technical indicators \cite{singh2017stock,nakano2018bitcoin}, technical and macro variables with denoising or representation learning \cite{bao2017deep,chong2017deep}, correlation-aware feature organization \cite{gunduz2017intraday}, and recurrent return sequences \cite{fischer2018deep,shen2018gru}. These outputs can be used as expected returns, risk states, or regime probabilities in a downstream portfolio or trading problem, where the quantum part is not just a vague metaphor but a QUBO or annealing-style decision module.

Sakuler et al. provide a useful recent benchmark because their 2025 study tests portfolio optimization on a ``real-world problem'' where a classical solution is already used in production \cite{sakuler2025real}. They formulate Markowitz-style portfolio optimization as QUBO, solve it with D-Wave QBSolv, Hybrid BQM, and Hybrid CQM, and benchmark the results against an exact classical strategy \cite{sakuler2025real}. Their conclusion is balanced: hybrid quantum approaches can find solutions ``close to the exact optimum'' in return and volatility, but QUBO parameter tuning and constraint handling are crucial; in the larger 499-asset experiment, Hybrid CQM gives more consistent feasible portfolios than QBSolv, while some QUBO samples violate the volatility constraint and must be discarded \cite{sakuler2025real}. This is directly relevant to our project because it shows both the application value and the implementation difficulty of quantum finance.

This integration still needs to be tested financially, not only statistically. Fischer and Krauss show that predictive strength may weaken after transaction costs and regime changes \cite{fischer2018deep}, while Nakano et al. show that bid--ask spread assumptions can change whether an ANN strategy beats buy-and-hold \cite{nakano2018bitcoin}. The quantum-finance layer should therefore be compared with a non-quantum baseline using return, volatility, constraint satisfaction, transaction costs, and risk-adjusted performance.

Herman et al. also make clear why this comparison matters: quantum Monte Carlo may offer quadratic speedups for risk modeling and derivative pricing, but only under resource assumptions that are not yet routine in practice, and quantum-native ML still lacks clear advantage on classical finance data \cite{herman2023quantum}. In other words, the literature supports a cautious hybrid framing rather than a claim that quantum computing has already solved financial engineering problems.

In general, the literature papers have supported a foundational layered design: Artificial Neural Networks provide the market-prediction component, while quantum finance contributes optimization, simulation, and risk-analysis tools. Their value of this integration should be judged by whether it improves feasible portfolio construction and financial risk control, not simply by whether the model uses quantum terminology.


\bibliographystyle{IEEEtran}
\bibliography{references}

\end{document}
